\chapter{Будущая работа и выбор узкой задачи}

\section{Рабочая тема}
Предлагаемое направление по умолчанию:
\begin{quote}
\textbf{Геометрические критерии локальной связности и проверяемые конечные
аппроксимации в комплексной динамике.}
\end{quote}
После пилота зафиксировать класс объектов, формат сертификата и основную
теорему. Обсудить постановку с руководителем.

\section{Три варианта специализации}
\begin{description}
\item[A. Геометрия и сертифицированные приближения --- основной кандидат.]
Построить конечные сертификаты включения в компоненту для заданного класса
геометрических конфигураций. Получить оценки размеров, глубины и точности
через геометрические запасы.
\item[B. Алгебраическая динамика конечного посткритического типа.]
Выбрать одно утверждение Bartholdi--Dudko о заданном классе разветвлённых
накрытий: алгоритм с доказательством корректности или структурную теорему.
\item[C. Аналитическая ренормализация.]
Восстановить один геометрический переход из principal nest или a priori bounds,
затем сформулировать и проверить усиление.
\end{description}
После 36-й недели оставляется одна основная ветка. Вторая может занимать
не более одной короткой сессии в две недели; третья откладывается.

\section{Что именно исследовать в ветке A}
Текущая схема использует
\[
O_N=\{c:|c|\le2,\ |f_c^n(0)|\le2\ (n\le N)\},\qquad
\bigcap_N O_N=\M.
\]
Для \(x\in\M\) обозначим
\[
C_N(x,r)=\Comp_x\bigl(O_N\cap\overline B(x,r)\bigr).
\]
Нужное достаточное условие имеет порядок кванторов
\begin{equation}\label{eq:buffer}
\forall x\in\M\ \forall\varepsilon>0\ \exists\,0<\delta<r<\varepsilon\quad
\forall N\ \exists L\ge N:
O_L\cap\overline B(x,\delta)\subseteq C_N(x,r).
\end{equation}
Радиусы выбираются \emph{до} \(N\), а \(L\) может от него зависеть.
Достаточность этого условия доказана в
\repo{Mlc/ParameterComponentApproximation.lean}.

В \repo{Mlc/FiniteComponentCriterion.lean} поле
\repo{LocalPiece.control} содержит \(\forall N\,\exists L\) и включение
в компоненту. Задача пилота --- задать конечные геометрические данные,
из которых следует такое включение в выбранном классе.

\textbf{Спецификация пилота.}
\begin{quote}
Для явно определённого класса геометрических конфигураций \(\mathcal C\)
построить формат рационального сертификата \(W\), проверяющий включение
\(O_L\cap\overline B(x,\delta)\subseteq C_N(x,r)\), и доказать:
(i) корректность проверки; (ii) существование сертификата при указанных
геометрических запасах; (iii) явную зависимость требуемой глубины и точности
от этих запасов.
\end{quote}
Начать с фиксированных \(x,N,r,\delta,L\). Переход к общим
кванторам~\eqref{eq:buffer} выделить в отдельную теорему.

Если использовать граф клеток, нужно доказать связность самих клеток,
достоверность их пересечений и нахождение в нужном внешнем множестве.
Для полуалгебраической постановки конечного уровня рассмотреть точные
методы реальной алгебраической геометрии и оценить их сложность.

\section{Проверка силы достаточного условия}
Локальная связность предела не влечёт буферное свойство произвольных
компактных внешних приближений. Контрпример:
\[
S=\{0\}\subset\R,\qquad
O_N=\{0\}\cup\{1/k:k\ge N+1\}.
\]
Все \(O_N\) компактны, убывают и имеют пересечение \(S\), которое локально
связно. Однако компонента нуля в \(O_N\cap[-r,r]\) равна \(\{0\}\):
подмножества прямой, содержащие два различных числа и связные, должны
содержать интервал, а этот набор счётен. Для любых \(L\) и \(\delta>0\)
в \(O_L\cap[-\delta,\delta]\) есть положительное \(1/k\). Значит, нужное
буферное включение не выполняется.

\section{Проверка новизны}
На 24-й и 36-й неделях для двух кандидатов составить таблицу: точная
формулировка; ближайшая теорема; что изменено; почему это не непосредственное
следствие; кто может проверить; минимальный контрпример к чрезмерному усилению.
Проверить доказательства и библиографии ближайших работ, затем передать
сравнение внешнему рецензенту.

\section{Каркас будущей рукописи}
\begin{enumerate}
\item Задача, точная общность и отличие от известных результатов.
\item Необходимые аналитические и топологические основания с доказательствами.
\item Ограничения прежних моделей; авторство и история каждого результата.
\item Основная новая конструкция или теорема.
\item Примеры, границы применимости и независимые проверки.
\item Lean-реализация: соответствие математики определениям, доверенная база,
воспроизводимость; оставшиеся обязательства вынесены явно.
\end{enumerate}
Обзор из \repo{docs/} использовать как справочник при подготовке рукописи.
